\chapter{Psoriasis}

\section*{Definition}
Psoriasis is a chronic, immune-mediated inflammatory skin disease characterized by well-demarcated, erythematous plaques with silvery-white scale, affecting 2--3\% of the global population. Plaque psoriasis (psoriasis vulgaris) accounts for 80--90\% of cases.

\section*{What Happens in Your Body}
Psoriasis involves a dysregulated IL-23/Th17 immune axis with excessive production of IL-17, IL-22, and TNF-\ensuremath{\alpha} by activated T-cells and innate lymphoid cells. These immune signaling proteins drive hyperproliferation of keratinocytes (epidermal turnover from 28 days to 3--5 days), aberrant differentiation, and angiogenesis. The Koebner phenomenon (lesions at trauma sites) reflects this immune hyperresponsiveness.

\section*{Causes}
\begin{itemize}
  \item \textbf{Genetic:} HLA-C*06:02 (strongest risk allele), IL23R, IL12B, TNIP1, and other susceptibility loci; 40\% have family history
  \item \textbf{Environmental triggers:} Streptococcal pharyngitis (guttate psoriasis), stress, trauma (Koebner), cold weather
  \item \textbf{Drug-induced:} Beta-blockers, lithium, antimalarials, NSAIDs, ACE inhibitors, TNF-\ensuremath{\alpha} inhibitors (paradoxical), interferon
  \item \textbf{Comorbidities:} Psoriatic arthritis (30\%), metabolic syndrome, cardiovascular disease, depression, inflammatory bowel disease
\end{itemize}

\section*{When to See a Doctor}
See your doctor if:
\begin{itemize}
  \item Widespread plaques, severe itching, or pain interfering with quality of life
  \item Joint pain, morning stiffness, dactylitis, or enthesitis suggesting psoriatic arthritis
  \item Erythrodermic psoriasis (\ensuremath{>}90\% body surface area) or generalized pustular psoriasis --- medical emergencies
  \item Scalp, nail, or genital involvement that impairs function
\end{itemize}

\section*{Related Symptoms}
\begin{itemize}
  \item Psoriatic arthritis, nail pitting, onycholysis, guttate psoriasis, inverse psoriasis, pustular psoriasis
\end{itemize}
\textbf{Important:} Nail psoriasis (pitting, onycholysis, subungual hyperkeratosis) may occur in isolation or precede cutaneous lesions and is a significant predictor of concurrent psoriatic arthritis.

\section*{Self-Care / Home Management}
\begin{itemize}
  \item Daily emollient application and gentle scale removal with moisturizing cleansers
  \item Moderate sun exposure (10--15 minutes daily) for vitamin D effect; avoid sunburn
  \item Stress management, smoking cessation, and alcohol moderation to reduce flares
\end{itemize}

\section*{How Your Doctor Will Evaluate You}
\begin{itemize}
  \item Clinical diagnosis: Sharply demarcated erythematous plaques with micaceous scale on extensor surfaces, scalp, umbilicus, and gluteal cleft
  \item Skin biopsy with histology showing acanthosis, parakeratosis, hypogranulosis, dilated capillaries, and Munro microabscesses
  \item Screening for psoriatic arthritis: PEST or PASE questionnaire; periodic C-reactive protein and joint examination
\end{itemize}

\section*{Treatment Options}
\begin{itemize}
  \item Mild--moderate (<5\% BSA): Topical corticosteroids (class 1--4), vitamin D analogs (calcipotriol), tazarotene, or fixed-combination products
  \item Moderate--severe: Phototherapy (narrowband UVB, PUVA), oral systemic agents (methotrexate, cyclosporine, apremilast)
  \item Biologics: TNF-\ensuremath{\alpha} inhibitors (adalimumab, etanercept), IL-17 inhibitors (secukinumab, ixekizumab), IL-23 inhibitors (guselkumab, risankizumab)
  \item Treat comorbidities: Screen for metabolic syndrome, depression, and cardiovascular risk factors
\end{itemize}

\section*{References}
\begin{thebibliography}{99}
\bibitem{psoriasis1} Lebwohl MG. "Psoriasis." Ann Intern Med. 2018;168(7):ITC49-ITC64.
\bibitem{psoriasis2} Griffiths CEM, Armstrong AW, Gudjonsson JE, et al. "Psoriasis." Lancet. 2021;397(10281):1301-1315.
\bibitem{psoriasis3} Menter A, Gelfand JM, Connor C, et al. "Joint American Academy of Dermatology-National Psoriasis Foundation guidelines of care for the management of psoriasis with systemic nonbiologic therapies." J Am Acad Dermatol. 2020;82(6):1445-1486.
\bibitem{psoriasis4} Menter A, Korman NJ, Elmets CA, et al. "Guidelines of care for the management of psoriasis and psoriatic arthritis." J Am Acad Dermatol. 2009;60(4):643-659.
\bibitem{psoriasis5} Armstrong AW, Read C. "Pathophysiology, clinical presentation, and treatment of psoriasis: a review." JAMA. 2020;323(19):1945-1960.
\bibitem{psoriasis6} Boehncke WH, Schon MP. "Psoriasis." Lancet. 2015;386(9997):983-994.
\end{thebibliography}

\clearpage
